\section{Open Points and Next Steps}
\label{sec:next}

\begin{enumerate}[leftmargin=*,itemsep=3pt]
\item \textbf{Fluent verification (CFD team).} Run the selected geometry in Fluent (density-based,
k-$\omega$ SST, ideal gas, Sutherland viscosity) at $M = 0.95$, 1.2 and 1.47, plus the
$5^\circ$ incidence case (R3). This replaces the transonic assumption of
Section~\ref{sec:transonic} and checks the Euler wave drag with viscous effects; the Euler
wall-pressure files are the reference to compare against. The setup is detailed in
\texttt{CFD/Fluent\_Workflow\_Nose\_Cone.pdf} of the repository (Appendix~\ref{app:scripts}).
\item \textbf{RocketPy update (simulation team).} Replace the LVS nose in the RocketPy model by
the final geometry and its $C_D(M)$, then re-check apogee, the \SI{30}{m/s} rail exit speed and
the static margin (2--4 cal, 2027\_C\_SE\_\allowbreak FD\_REQ\_07) once the fins are defined.
\item \textbf{Requirements to create} (none exists for the Weisshorn nosecone). Proposed:
\begin{itemize}[itemsep=1pt]
\item profile: Von K\'arm\'an, $L = \FINL{}$ ($L/D=4.5$, \FINLPHYS{} physical with the tip), coordinates from \texttt{nose\_profile\_final.csv}, radius tolerance $\pm\SI{0.5}{mm}$;
\item nosecone bay length (curved part + the cylinder that completes the 3U housing) \FINBAY{}, to be booked in the \SI{5.4}{m} LV budget (2027\_C\_SE\_REQ\_04);
\item tip: spherical, radius \SI{4.5}{mm}.
\end{itemize}
\end{enumerate}
